\paragraph{Chapter Structure}

\begin{itemize}
\item In Chapter~{\ref{sec:psRB-language}}, we first give the language and program model formally similar to our language mode in
\redd{PART \romannum{1}}.

\item  Chapter~\ref{sec:psRB-definition} gives the \emph{reachability-bound} definition.

\item Chapter~\ref{sec:psRB-analysis-alg} presents the new \emph{path-sensitive reachability-bound} analysis algorithm, named $\kw{psRB}$ through the following sections.
%
Section~\ref{sec:psRB-preface} first gives the main steps of this algorithm as a structural section. There are five steps in this algorithm.
In Section~\ref{sec:psRB-progabs}, the first step abstracted the program into an \emph{abstract transition graph}.
Based on this graph, Section~\ref{sec:psRB-refine} then refines the multi-paths loop in this program and generates a refined program.
According to the graph and the refined program, Section~\ref{sec:psRB-rank} computes the ranking function for each edge on the graph and estimates the value of each ranking function.
Section~\ref{sec:psRB-psrb} gives the \emph{path-sensitive reachability-bound} algorithm and the rest of the sections give the computation algorithm for each step. Specifically, 
Section~\ref{sec:psRB-pathlocalrb} gives the local \emph{path reachability-bound} estimation algorithm, Section~\ref{sec:psRB-looprb} presents the \emph{loop reachability-bound} estimation algorithm,
and Section~\ref{sec:psRB-pathrb} gives the global  \emph{path reachability-bound} estimation algorithm.

\item  {Chapter~{\ref{sec:psRB-eval}}} presents four step-by-step examples to illustrate the algorithm clearly in Section~\ref{sec:psRB-example} firstly,
and then shows the implementation and evaluation results over five real world benchmarks in Section~\ref{sec:psRB-eval}.

\item {Chapter~\ref{sec:psRB-relatedwork}} discusses the related works in this area.
\end{itemize}

\paragraph{Contributions}
The contributions of this work can be summarized as follows,
\begin{itemize}
  \item We introduce a path-sensitive reachability-bound algorithm 
  combining \emph{amortized bound analysis}  and \emph{path refinement}. The algorithm computes tight bounds on the times that each program point is evaluated.

  \item We identify two quantities,  \emph{path reachability-bound} and \emph{loop reachability-bound}, which are useful to measure precise information about the program execution.  We also propose two algorithms for soundly estimating these two quantities.
  \item We develop a prototype implementation with evaluation over 290 programs.
  The evaluation shows improvement in precision compared to other tools for reachability-bound analysis and the worst-case complexity analysis.

\end{itemize}